\documentclass{article}

\usepackage[margin=1in]{geometry}
\linespread{1.2}
\usepackage{amsmath,amsthm,amssymb,amsfonts,commath}


\newcommand{\mto}{\stackrel{\mu}{\rightarrow}}
\newcommand{\Mcal}{\mathcal{M}}

\newcommand{\TV}{\mathrm{TV}}
\newcommand{\AC}{\mathrm{AC}}

 
\title{Real Analysis I --- Homework 6}

% Replace with your name and ID below
\author{FirstName LastName \\ ID: 123-45-6789}

\date{}


\begin{document}

\maketitle

Unless otherwise specified, all functions are defined on $E \subset \mathbb{R}^{n}$ where $E \in \mathcal{M}$.
\bigskip

%For notation simplicity, a set such as $\{x \in E: f(x) > t\}$ is written as $\{f>t\}$ for short.

\begin{enumerate}

\item If $\nu$ is a signed measure, then $E$ is $\nu$-null iff $|\nu|(E) = 0$. Also, if $\nu$ and $\mu$ are signed measures, $\nu \perp \mu$ iff $|\nu| \perp \mu$ iff $\nu^{+} \perp \mu$ and $\nu^{-} \perp \mu$.

%uncomment below and write your proof after it.
%\begin{proof}
%\end{proof}


%\newpage % uncomment this line to start on a new page

\item Suppose $\{\nu_{j}\}$ is a sequence of positive measures. If $\nu_{j} \perp \mu$ for all $j$, then $(\sum_{j=1}^{\infty} \nu_{j}) \perp \mu$; and if $\nu_{j} \ll \mu$ for all $j$, then $(\sum_{j=1}^{\infty}  \nu_{j}) \ll \mu$.

%uncomment below and write your proof after it.
%\begin{proof}
%\end{proof}


%\newpage % uncomment this line to start on a new page

\item Suppose that $\mu, \nu$ are positive measures on $(X, \Mcal)$ with $\nu \ll \mu$, and let $\lambda = \mu + \nu$. If $f = \frac{\dif \nu}{\dif \lambda}$, then $0 \le f < 1$ $\mu$-a.e., and $\frac{\dif \nu}{\dif \mu} = \frac{f}{1-f}$.

%uncomment below and write your proof after it.
%\begin{proof}
%\end{proof}



%\newpage % uncomment this line to start on a new page

\item Suppose $f \in \mathrm{BV}([a,b])$, and $f(x) \ge c > 0$ on $[a,b]$. Show that $\frac{1}{f} \in \mathrm{BV}([a,b])$.

%uncomment below and write your proof after it.
%\begin{proof}
%\end{proof}


%\newpage % uncomment this line to start on a new page

\item Show that $f \in \mathrm{BV}([a,b])$ iff there exists a non-decreasing function $\phi: [a,b] \to \mathbb{R}$ such that
\[
|f(y) - f(x) | \le \phi(y) - \phi(x), \quad \mbox{for all}\ a \le x < y \le b.
\]

%uncomment below and write your proof after it.
%\begin{proof}
%\end{proof}



%\newpage % uncomment this line to start on a new page

\item Suppose $f \in \mathrm{AC}([a,b])$. Show that
\[
\int_{a}^{b} |f'(x)| \dif x = \mathrm{TV}(f, [a,b]).
\]
Hint: We showed $\ge$ in class. For $\le$, it is easy to show that $\int_{a}^{b} s(x) f'(x) \dif x \le \mathrm{TV}(f)$ where $s(x) = \sum_{i=1}^{p}c_{i}\chi_{[x_{i-1},x_{i})}(x)$ and $|c_{i}| \le 1$ for any partition $\Delta: a=x_{0} < x_{1} < \dots < x_{n} =b$. Let $\phi_{k}$ be a sequence of simple functions such that $\phi_{k} \to f'$, and $s_{k}(x)$ be the sign of $\phi_{k}(x)$. Apply DCT to $s_{k} f'$.

%uncomment below and write your proof after it.
%\begin{proof}
%\end{proof}


%\newpage % uncomment this line to start on a new page


\item Suppose $f \in L([a,b])$. Show that for a.e.~$x \in [a,b]$, there is
\[
\lim_{h \to 0} \frac{1}{h} \int_{x - h_{1}}^{x + h_{2}} |f(y) - f(x) | \dif y  = 0,
\]
where $h_{1},h_{2} \ge 0$ and $h = h_{1}+h_{2}$.

%uncomment below and write your proof after it.
%\begin{proof}
%\end{proof}




%\newpage % uncomment this line to start on a new page


\item Suppose $f \in L^{2}((0,\pi))$. Show that the following two inequalities cannot hold simultaneously:
\[
\int_{0}^{\pi}(f(x) - \sin x)^{2} \dif x \le \frac{4}{9}, \quad 
\int_{0}^{\pi}(f(x) - \cos x)^{2} \dif x \le \frac{1}{9}.
\]

%uncomment below and write your proof after it.
%\begin{proof}
%\end{proof}



%\newpage % uncomment this line to start on a new page


\item Suppose $p,q,r > 0$ satisfy $\frac{1}{p} + \frac{1}{q} + \frac{1}{r} = 1$. Then for any functions $f \in L^{p}(E)$, $g \in L^{q}(E)$ and $h \in L^{r}(E)$, there is
\[
\int_{E} |f g h| \dif x \le \|f\|_{p} \|g\|_{q} \|h\|_{r}.
\]

%uncomment below and write your proof after it.
%\begin{proof}
%\end{proof}




%\newpage % uncomment this line to start on a new page


\item Suppose $f \in L^{p}(\mathbb{R})$ where $1 < p < \infty$, and $q$ is the conjugate of $p$. Let $F(x) = \int_{a}^{x} f(t) \dif t$. Show that $F(x+h) - F(x) = o(|h|^{1/q})$ as $h \to 0$ for all $x$.

%uncomment below and write your proof after it.
%\begin{proof}
%\end{proof}





\end{enumerate}


\end{document}